% GENERATED FILE. Edit canonical lessons, then run npm run generate:books.
% canonical-source-hash: fnv1a64:decfdde8be16c5ea
% canonical-lessons: ML-C543-bonasu, ML-C543-apeksikkuka, ML-C543-jolikketukkuka, ML-C543-piriccuvituka, ML-C543-tolilutama

\chapter{Hired and Fired}
\label{ch:ml-hired-and-fired}
\hlchaptermodality{543}

\begin{chapteropening}
\textbf{By the end of this chapter:} \emph{I can say a bonus, to apply, to take on staff, to dismiss and an employer.}

Complete the last lesson of chapter 543: I can say a bonus, to apply, to take on staff, to dismiss and an employer.
\end{chapteropening}

\section[\texorpdfstring{\ml{ബോണസ്} (bōṇasŭ)}{ബോണസ് (bōṇasŭ)}]{\ml{ബോണസ്} (bōṇasŭ) — a bonus}

\label{lesson:ML-C543-bonasu}



\begin{quote}
Before the new one: say the Malayalam for pay, then the Malayalam for to step down.
\end{quote}

\subsection*{You'll want to know: \ml{ബോണസ്}}
\textbf{\ml{ബോണസ്}} — \emph{bōṇasŭ} — \textquotedblleft{}a bonus\textquotedblright{}.

\begin{quote}
\ml{വർഷാവസാനം} \ml{ജീവനക്കാർക്ക്} \ml{ബോണസ്} \ml{കിട്ടി}.
\end{quote}

\begin{quote}
\emph{varṣāvasānaṁ jīvanakkārkkŭ bōṇasŭ kiṭṭi.}
\end{quote}

\begin{quote}
At the end of the year the staff got a bonus.
\end{quote}

\ml{ബോണസ്} \ml{കിട്ടുക}, to get a bonus; it is paid on top of the salary.

\begin{cousinweb}
\textbf{\ml{ബോണസ്}} (\emph{bōṇasŭ}) was borrowed from English \emph{bonus}.
\end{cousinweb}

\begin{grammarlens}[title={What you've built}]
One more word for describing what you lived and felt.
\end{grammarlens}

\subsection*{Your turn}
\begin{itemize}
\raggedright
  \item \emph{Say it:} \emph{bōṇasŭ}
  \item \emph{Say it:} \emph{bōṇasŭ}, once more
  \item \emph{Say it:} \emph{varṣāvasānaṁ jīvanakkārkkŭ bōṇasŭ kiṭṭi.} — At the end of the year the staff got a bonus.
  \item \emph{Recall:} say the Malayalam for pay, then the Malayalam for to step down, then say \emph{bōṇasŭ} again
\end{itemize}

\begin{tcolorbox}[breakable,colback=teal!4,colframe=teal!35!black,title={Before you move on}]
What does \ml{ബോണസ്} mean? (\textquotedblleft{}A bonus\textquotedblright{}.) Say it once more.
\end{tcolorbox}
\section[\texorpdfstring{\ml{അപേക്ഷിക്കുക} (apēkṣikkuka)}{അപേക്ഷിക്കുക (apēkṣikkuka)}]{\ml{അപേക്ഷിക്കുക} (apēkṣikkuka) — to apply; to request}

\label{lesson:ML-C543-apeksikkuka}



\begin{quote}
Before the new one: say the Malayalam for to step down, then the Malayalam for a bonus.
\end{quote}

\subsection*{You'll want to know: \ml{അപേക്ഷിക്കുക}}
\textbf{\ml{അപേക്ഷിക്കുക}} — \emph{apēkṣikkuka} — \textquotedblleft{}to apply; to request\textquotedblright{}.

\begin{quote}
\ml{പുതിയ} \ml{ജോലിക്ക്} \ml{ഞാൻ} \ml{ഓൺലൈനായി} \ml{അപേക്ഷിച്ചു}.
\end{quote}

\begin{quote}
\emph{putiya jōlikkŭ ñān ōṇlaināyi apēkṣiccu.}
\end{quote}

\begin{quote}
I applied online for a new job.
\end{quote}

Past: \ml{അപേക്ഷിച്ചു}, \emph{apēkṣiccu}.

The noun \ml{അപേക്ഷ}, a request or an application, is already familiar; what one applies for takes -\ml{ക്ക്}.

\begin{cousinweb}
\textbf{\ml{അപേക്ഷിക്കുക}} (\emph{apēkṣikkuka}) is built from \ml{അപേക്ഷ}, from Sanskrit \emph{apekṣā}, \textquotedblleft{}looking about, regard, expectation\textquotedblright{} (Monier-Williams), with the Malayalam verb ending -\ml{ിക്കുക}.
\end{cousinweb}

\begin{grammarlens}[title={What you've built}]
One more word for describing what you lived and felt.
\end{grammarlens}

\subsection*{Your turn}
\begin{itemize}
\raggedright
  \item \emph{Say it:} \emph{apēkṣikkuka}
  \item \emph{Say it:} \emph{apēkṣikkuka}, once more
  \item \emph{Say it:} \emph{putiya jōlikkŭ ñān ōṇlaināyi apēkṣiccu.} — I applied online for a new job.
  \item \emph{Recall:} say the Malayalam for to step down, then the Malayalam for a bonus, then say \emph{apēkṣikkuka} again
\end{itemize}

\begin{tcolorbox}[breakable,colback=teal!4,colframe=teal!35!black,title={Before you move on}]
What does \ml{അപേക്ഷിക്കുക} mean? (\textquotedblleft{}To apply; to request\textquotedblright{}.) Say it once more.
\end{tcolorbox}
\section[\texorpdfstring{\ml{ജോലിക്കെടുക്കുക} (jōlikkeṭukkuka)}{ജോലിക്കെടുക്കുക (jōlikkeṭukkuka)}]{\ml{ജോലിക്കെടുക്കുക} (jōlikkeṭukkuka) — to take on staff; to hire}

\label{lesson:ML-C543-jolikketukkuka}



\begin{quote}
Before the new one: say the Malayalam for a bonus, then the Malayalam for to apply.
\end{quote}

\subsection*{You'll want to know: \ml{ജോലിക്കെടുക്കുക}}
\textbf{\ml{ജോലിക്കെടുക്കുക}} — \emph{jōlikkeṭukkuka} — \textquotedblleft{}to take on staff; to hire\textquotedblright{}.

\begin{quote}
\ml{ആ} \ml{കട} \ml{രണ്ട്} \ml{വിദ്യാർത്ഥികളെ} \ml{ജോലിക്കെടുത്തു}.
\end{quote}

\begin{quote}
\emph{ā kaṭa raṇṭŭ vidyārtthikaḷe jōlikkeṭuttu.}
\end{quote}

\begin{quote}
That shop took on two students.
\end{quote}

Past: \ml{ജോലിക്കെടുത്തു}, \emph{jōlikkeṭuttu}.

\ml{ജോലിക്ക്} for work + \ml{എടുക്കുക} to take.

\begin{cousinweb}
\textbf{\ml{ജോലിക്കെടുക്കുക}} (\emph{jōlikkeṭukkuka}) is a Malayalam compound built from \ml{ജോലിക്ക്}, \textquotedblleft{}for work\textquotedblright{}, and \ml{എടുക്കുക}, \textquotedblleft{}to take up\textquotedblright{}, which is native (DEDR 851).
\end{cousinweb}

\begin{grammarlens}[title={What you've built}]
One more word for describing what you lived and felt.
\end{grammarlens}

\subsection*{Your turn}
\begin{itemize}
\raggedright
  \item \emph{Say it:} \emph{jōlikkeṭukkuka}
  \item \emph{Say it:} \emph{jōlikkeṭukkuka}, once more
  \item \emph{Say it:} \emph{ā kaṭa raṇṭŭ vidyārtthikaḷe jōlikkeṭuttu.} — That shop took on two students.
  \item \emph{Recall:} say the Malayalam for a bonus, then the Malayalam for to apply, then say \emph{jōlikkeṭukkuka} again
\end{itemize}

\begin{tcolorbox}[breakable,colback=teal!4,colframe=teal!35!black,title={Before you move on}]
What does \ml{ജോലിക്കെടുക്കുക} mean? (\textquotedblleft{}To take on staff; to hire\textquotedblright{}.) Say it once more.
\end{tcolorbox}
\section[\texorpdfstring{\ml{പിരിച്ചുവിടുക} (piriccuviṭuka)}{പിരിച്ചുവിടുക (piriccuviṭuka)}]{\ml{പിരിച്ചുവിടുക} (piriccuviṭuka) — to dismiss; to sack}

\label{lesson:ML-C543-piriccuvituka}



\begin{quote}
Before the new one: say the Malayalam for to apply, then the Malayalam for to take on staff.
\end{quote}

\subsection*{You'll want to know: \ml{പിരിച്ചുവിടുക}}
\textbf{\ml{പിരിച്ചുവിടുക}} — \emph{piriccuviṭuka} — \textquotedblleft{}to dismiss; to sack\textquotedblright{}.

\begin{quote}
\ml{നഷ്ടം} \ml{കാരണം} \ml{കമ്പനി} \ml{നൂറ്} \ml{പേരെ} \ml{പിരിച്ചുവിട്ടു}.
\end{quote}

\begin{quote}
\emph{naṣṭaṁ kāraṇaṁ kampani nūṟŭ pēre piriccuviṭṭu.}
\end{quote}

\begin{quote}
Because of losses the company laid off a hundred people.
\end{quote}

Past: \ml{പിരിച്ചുവിട്ടു}, \emph{piriccuviṭṭu}.

\ml{പിരിച്ചു}, having separated, + \ml{വിടുക}, to let go.

\begin{cousinweb}
\textbf{\ml{പിരിച്ചുവിടുക}} (\emph{piriccuviṭuka}) is native Dravidian, a compound of \ml{പിരിക്കുക}, \textquotedblleft{}to sever, to dismiss\textquotedblright{} (DEDR 4176), and \ml{വിടുക}, \textquotedblleft{}to let go\textquotedblright{} (DEDR 5393).
\end{cousinweb}

\begin{grammarlens}[title={What you've built}]
One more word for describing what you lived and felt.
\end{grammarlens}

\subsection*{Your turn}
\begin{itemize}
\raggedright
  \item \emph{Say it:} \emph{piriccuviṭuka}
  \item \emph{Say it:} \emph{piriccuviṭuka}, once more
  \item \emph{Say it:} \emph{naṣṭaṁ kāraṇaṁ kampani nūṟŭ pēre piriccuviṭṭu.} — Because of losses the company laid off a hundred people.
  \item \emph{Recall:} say the Malayalam for to apply, then the Malayalam for to take on staff, then say \emph{piriccuviṭuka} again
\end{itemize}

\begin{tcolorbox}[breakable,colback=teal!4,colframe=teal!35!black,title={Before you move on}]
What does \ml{പിരിച്ചുവിടുക} mean? (\textquotedblleft{}To dismiss; to sack\textquotedblright{}.) Say it once more.
\end{tcolorbox}
\section[\texorpdfstring{\ml{തൊഴിലുടമ} (toḻiluṭama)}{തൊഴിലുടമ (toḻiluṭama)}]{\ml{തൊഴിലുടമ} (toḻiluṭama) — an employer}

\label{lesson:ML-C543-tolilutama}



\begin{quote}
Before the new one: say the Malayalam for to take on staff, then the Malayalam for to dismiss.
\end{quote}

\subsection*{You'll want to know: \ml{തൊഴിലുടമ}}
\textbf{\ml{തൊഴിലുടമ}} — \emph{toḻiluṭama} — \textquotedblleft{}an employer\textquotedblright{}.

\begin{quote}
\ml{തൊഴിലുടമ} \ml{ശമ്പളം} \ml{കൃത്യസമയത്ത്} \ml{കൊടുത്തു}.
\end{quote}

\begin{quote}
\emph{toḻiluṭama śampaḷaṁ kṛtyasamayattŭ koṭuttu.}
\end{quote}

\begin{quote}
The employer paid the wages on time.
\end{quote}

\ml{തൊഴിൽ} work + \ml{ഉടമ} an owner.

\begin{cousinweb}
\textbf{\ml{തൊഴിലുടമ}} (\emph{toḻiluṭama}) is native Dravidian, a compound of \ml{തൊഴിൽ}, \textquotedblleft{}an occupation\textquotedblright{} (DEDR 3524), and \ml{ഉടമ}, \textquotedblleft{}property, an owner\textquotedblright{} (DEDR 593).
\end{cousinweb}

\begin{grammarlens}[title={What you've built}]
One more word for describing what you lived and felt.
\end{grammarlens}

\subsection*{Your turn}
\begin{itemize}
\raggedright
  \item \emph{Say it:} \emph{toḻiluṭama}
  \item \emph{Say it:} \emph{toḻiluṭama}, once more
  \item \emph{Say it:} \emph{toḻiluṭama śampaḷaṁ kṛtyasamayattŭ koṭuttu.} — The employer paid the wages on time.
  \item \emph{Recall:} say the Malayalam for to take on staff, then the Malayalam for to dismiss, then say \emph{toḻiluṭama} again
\end{itemize}

\begin{tcolorbox}[breakable,colback=teal!4,colframe=teal!35!black,title={Before you move on}]
What does \ml{തൊഴിലുടമ} mean? (\textquotedblleft{}An employer\textquotedblright{}.) Say it once more.
\end{tcolorbox}
